% =============================================================================
% 1.9 Correspondance approximative, résolution d'entités et graphes
% =============================================================================
\section{Résolution d'entités et graphes de connaissances}
\label{sec:ch1-entites}

\subsection{Correspondance approximative}

La correspondance approximative (\anglais{fuzzy matching}) désigne l'ensemble
des techniques permettant de comparer des chaînes malgré des variations
orthographiques, des fautes de frappe ou des artefacts d'OCR. Elle
s'inscrit dans un héritage ancien --- la distance d'édition de Levenshtein
(1965), les codes phonétiques de type Soundex, les coefficients de Jaccard
et les $n$-grammes --- et demeure indispensable dès que l'on traite des
textes produits par des humains ou par une reconnaissance optique imparfaite.

\begin{btptable}{Méthodes fondamentales de correspondance approximative}
\label{tab:fuzzy}
\begin{tabularx}{\linewidth}{@{}L{2.8cm} Y L{3.6cm}@{}}
\toprule
\rowcolor{btppale}
\btpth{Méthode} & \btpth{Principe} & \btpth{Application typique} \\
\midrule
Levenshtein & Nombre d'insertions, suppressions et substitutions. & Fautes, artefacts d'OCR. \\
Jaccard & Intersection sur union d'ensembles de tokens. & Ordre des mots secondaire. \\
$n$-grammes & Partage de sous-séquences de $n$ caractères. & Texte long, OCR. \\
Soundex et variantes & Encodage phonétique. & Noms propres (limité en français). \\
\bottomrule
\end{tabularx}
\end{btptable}

\subsection{Résolution d'entités}

La résolution d'entités consiste à décider que deux libellés désignent la
même entité réelle. Elle s'appuie sur la correspondance approximative, mais
l'intègre dans un processus~: normalisation (casse, accents, formes
juridiques), comparaison, seuillage, et, en cas de doute, arbitrage humain.

La figure~\ref{fig:resolution-entites} en illustre le déroulement sur deux
dénominations d'une même entreprise.

\begin{figure}[htbp]
  \centering
  \resizebox{0.92\linewidth}{!}{% Chaîne de résolution d'entités
\begin{tikzpicture}[node distance=7mm and 14mm]

  \node[btpboxplain, text width=3.4cm] (a) {\ttfamily\scriptsize «~SITINFRA Sarl~»};
  \node[btpboxplain, text width=3.4cm, right=26mm of a] (b) {\ttfamily\scriptsize «~Sitinfra Cameroun~»};

  \node[btpbox, text width=3.4cm, below=of a] (na) {Normalisation\\[-1pt]\btplabel{casse, accents, formes juridiques}};
  \node[btpbox, text width=3.4cm, below=of b] (nb) {Normalisation\\[-1pt]\btplabel{casse, accents, formes juridiques}};

  \node[btpboxplain, text width=3.4cm, below=of na] (ra) {\ttfamily\scriptsize sitinfra};
  \node[btpboxplain, text width=3.4cm, below=of nb] (rb) {\ttfamily\scriptsize sitinfra cameroun};

  \node[btpboxalt, text width=5.2cm, below=10mm of {$(ra)!0.5!(rb)$}] (cmp)
    {Comparaison approximative\\[-1pt]\btplabel{Levenshtein, Jaccard, $n$-grammes}};

  \node[btpdecision, text width=2.1cm, below=9mm of cmp] (dec)
    {Score $>$ seuil~?};

  \node[btpboxstrong, text width=3.0cm, below left=9mm and 4mm of dec] (same)
    {Même entité\\[-1pt]\btplabel{fusion}};
  \node[btpbox, text width=3.0cm, below right=9mm and 4mm of dec] (diff)
    {Entités distinctes\\[-1pt]\btplabel{arbitrage humain}};

  \draw[btparrow] (a) -- (na);
  \draw[btparrow] (b) -- (nb);
  \draw[btparrow] (na) -- (ra);
  \draw[btparrow] (nb) -- (rb);
  \draw[btparrow] (ra.south) -- (cmp.north west);
  \draw[btparrow] (rb.south) -- (cmp.north east);
  \draw[btparrow] (cmp) -- (dec);
  \draw[btparrow] (dec) -- node[btplabel, left, xshift=-1mm] {oui} (same);
  \draw[btparrow] (dec) -- node[btplabel, right, xshift=1mm] {non} (diff);

\end{tikzpicture}
}
  \caption{Chaîne de résolution d'entités}
  \label{fig:resolution-entites}
\end{figure}

Dans un corpus d'appels d'offres, les entités récurrentes sont nombreuses
et leurs dénominations instables~:

\begin{itemize}
  \item \btpterm{Entreprises}~: « SITINFRA », « Sitinfra Cameroun »,
        « SITINFRA Sarl »~;
  \item \btpterm{Projets}~: un même ouvrage désigné par son lieu, son maître
        d'ouvrage ou son acronyme interne~;
  \item \btpterm{Certifications}~: « ISO 9001 » et « ISO 9001:2015 »~;
  \item \btpterm{Personnes}~: variations d'état civil, de titre et d'ordre
        des noms.
\end{itemize}

Sans résolution, le système crée des doublons, manque des rapprochements
et, le cas échéant, peuple un graphe de n\oe uds redondants. Avec une
résolution trop agressive, il fusionne des entités distinctes. Le seuil
et l'arbitrage humain ne sont donc pas des raffinements~: ils sont
constitutifs de la méthode.

\subsection{Graphes de connaissances}

Un graphe de connaissances représente les entités comme des n\oe uds et
les relations comme des arêtes typées. La figure~\ref{fig:graphe}
en donne une instance schématique pour le domaine des appels d'offres.

\begin{figure}[htbp]
  \centering
  \resizebox{\linewidth}{!}{% Principe d'un graphe de connaissances appliqué aux appels d'offres
\begin{tikzpicture}[node distance=16mm and 26mm,
                    every edge quotes/.style={btplabel, fill=white, inner sep=1.5pt}]

  \node[btpboxstrong, text width=2.2cm] (ent) {ENTREPRISE};
  \node[btpbox, text width=2.0cm, right=of ent] (proj) {PROJET};
  \node[btpbox, text width=2.2cm, right=of proj] (client) {CLIENT};

  \node[btpboxalt, text width=2.4cm, below=of ent] (cert) {CERTIFICATION};
  \node[btpboxalt, text width=2.2cm, below=of proj] (pers) {PERSONNE};
  \node[btpboxalt, text width=2.2cm, below=of client] (exig) {EXIGENCE};

  \draw[btparrow] (ent) -- node[btplabel, fill=white, inner sep=1.5pt, above]
       {\textsc{a\_realise}} (proj);
  \draw[btparrow] (proj) -- node[btplabel, fill=white, inner sep=1.5pt, above]
       {\textsc{pour}} (client);
  \draw[btparrow] (ent) -- node[btplabel, fill=white, inner sep=1.5pt, left]
       {\textsc{detient}} (cert);
  \draw[btparrow] (ent) -- node[btplabel, fill=white, inner sep=1.5pt, pos=0.32, above, sloped]
       {\textsc{emploie}} (pers);
  \draw[btparrow] (pers) -- node[btplabel, fill=white, inner sep=1.5pt, pos=0.5, above, sloped]
       {\textsc{a\_travaille\_sur}} (proj);
  \draw[btparrow] (cert) -- node[btplabel, fill=white, inner sep=1.5pt, below]
       {\textsc{satisfait}} (exig);
  \draw[btparrowlight] (exig) -- node[btplabel, fill=white, inner sep=1.5pt, right]
       {\textsc{issue\_de}} (client);

  \node[btplabel, below=7mm of pers, text width=10cm]
    {Une question comme « quels ingénieurs certifiés ont travaillé sur un projet
     comparable pour ce maître d'ouvrage~? » se résout par un parcours de
     chemins dans le graphe, là où une recherche textuelle échouerait.};

\end{tikzpicture}
}
  \caption{Principe d'un graphe de connaissances appliqué aux appels d'offres}
  \label{fig:graphe}
\end{figure}

Quatre familles de questions, mal servies par une recherche textuelle,
deviennent formulables~:

\begin{enumerate}
  \item \btpterm{Traçabilité des compétences}~: quels ingénieurs ont
        travaillé sur des projets comparables~?
  \item \btpterm{Conformité relationnelle}~: quelles certifications
        couvrent quelles exigences du règlement de consultation~?
  \item \btpterm{Analyse transverse}~: quels maîtres d'ouvrage ont déjà
        attribué des marchés à l'entreprise~?
  \item \btpterm{Capitalisation des relations}~: projets, clients,
        équipements et compétences deviennent des chemins navigables.
\end{enumerate}

L'extraction automatique de tels graphes à partir de textes non structurés
reste toutefois un défi ouvert, résumé au tableau~\ref{tab:defis-graphe}.

\begin{btptable}{Défis de l'extraction de graphes de connaissances}
\label{tab:defis-graphe}
\begin{tabularx}{\linewidth}{@{}L{3.4cm} Y Y@{}}
\toprule
\rowcolor{btppale}
\btpth{Défi} & \btpth{Description} & \btpth{Piste} \\
\midrule
Résolution d'entités & Une même entité sous des formes variées. & Normalisation et correspondance approximative. \\
Extraction de relations & Identifier le type et le sens d'un lien. & Modèle contraint, schéma prédéfini. \\
Évolution & Le graphe doit s'enrichir sans se corrompre. & Fusion incrémentale, versions. \\
Évaluation & Mesurer l'exactitude des triplets extraits. & Validation humaine par échantillon. \\
\bottomrule
\end{tabularx}
\end{btptable}

\begin{btpwarn}[Positionnement]
Le graphe de connaissances n'est pas un prérequis de la première version
d'un système de génération d'offres. Il devient pertinent lorsque les
questions relationnelles multi-sauts constituent un besoin avéré, et
lorsque le corpus est suffisamment propre pour que l'extraction ne
produise pas plus de bruit que de signal. Nous le traitons donc comme une
perspective de seconde phase, et non comme un fondement de l'architecture
initiale.
\end{btpwarn}
